\section{Directed trees and coherent-state spectral laws}
\label{sec:coherent-spectrum}

A projective coherent state carries considerably more structure than an
isotropic random vector in its ambient Hilbert space.  This distinction
survives when the number of samples and the Hilbert-space dimension grow
at the same rate.  We describe the resulting spectral law by a sum over
set partitions.  Its weights are directed spanning-tree counts.  The
noncrossing partitions familiar from free probability emerge precisely
when the complex dimension tends to infinity.

This is a spatial-kernel moment calculation, in the sense of the
Euclidean random-matrix framework of Bordenave~\cite{prob_bordenave2008}.
The coherent phase is retained: after rescaling it produces a directed,
rather than a symmetric, graph Laplacian.  Lowest-Landau-level random
impurity models provide a related one-complex-dimensional setting
\cite{prob_brezin1984}.  We give a direct proof on projective space,
including arbitrary fixed sampling density.

\subsection{The spectral law at fixed dimension and sampling density}

Fix a positive integer $q$ and write
\[
 D_m=\dim H^0(\mathbb P^q,\mathcal O(m))=\binom{m+q}{q}.
\]
Let $n=n_m$ satisfy $n_m/D_m\longrightarrow\rho\in(0,\infty)$.
Choose independent Fubini--Study distributed rays $z_1,\ldots,z_n$
in $\mathbb P^q$ and arbitrary unit lifts.  Their degree-$m$ coherent
Gram matrix and empirical spectral measure are
\begin{equation}
 B^{(m)}_{ij}=\langle z_i,z_j\rangle^m,
 \qquad
 \nu_m=\frac1n\sum_{j=1}^n\delta_{\lambda_j(B^{(m)})}.
 \label{eq:coherent-gram-definition}
\end{equation}
Changing the phases of the lifts gives a diagonal unitary conjugation.
In particular, the spectrum is intrinsic.  Moreover, $B^{(m)}\geq0$
and $B^{(m)}_{ii}=1$.

For a set partition $\pi$ of $\{1,\ldots,p\}$, identify the vertices
of the directed cycle $1\to2\to\cdots\to p\to1$ according to the
blocks of $\pi$.  Denote the resulting directed multigraph by $G_\pi$.
Parallel edges are retained.  Loops contribute nothing to its
outdegree Laplacian
\[
 L_{aa}=\sum_{b\ne a}n_{ab},\qquad L_{ab}=-n_{ab}\quad(a\ne b).
\]
The graph is connected and Eulerian.  Its rooted cofactor
$\tau(G_\pi)=\det L^{(o)}$ is the number of directed spanning
in-arborescences rooted at $o$.  It is a positive integer and is
independent of $o$.  The one-vertex cofactor is defined to be one.

\begin{theorem}[Coherent spectral law]
\label{thm:coherent-spectral-law}
There is a unique probability measure $\nu_{q,\rho}$ on $[0,\infty)$
whose positive integer moments are
\begin{equation}
 M_p(q,\rho)=
 \sum_{\pi\in\mathfrak P(p)}
       \rho^{|\pi|-1}\tau(G_\pi)^{-q}.
 \label{eq:coherent-tree-moments}
\end{equation}
For every fixed $p\geq1$,
\[
 \mathbb E\!\int x^p\,d\nu_m\longrightarrow M_p(q,\rho),
 \qquad
 \operatorname{Var}\!\left(\int x^p\,d\nu_m\right)=O_p(n_m^{-1}).
\]
Consequently $\nu_m$ converges weakly in probability to
$\nu_{q,\rho}$, with convergence in probability of every fixed
positive integer moment.  If $\sum_m n_m^{-1}<\infty$, these
convergences hold almost surely on any probability space carrying
the rows.  In particular this applies along all degrees when $q\geq2$.

The same assertions hold after conditioning on any fixed finite
collection of bounded empirical moments lying within distance
$\epsilon_m$ of their Fubini--Study expectations, provided
$n_m\epsilon_m^2\to\infty$.
\end{theorem}

The first four limiting moments are
\begin{align}
 M_1&=1,& M_2&=1+\rho,& M_3&=1+3\rho+\rho^2,\notag\\
 M_4&=1+(6+2^{-q})\rho+6\rho^2+\rho^3.
 \label{eq:coherent-first-moments}
\end{align}
At unit sampling density and complex dimension two they are
$1,2,5,57/4$.  The sole crossing partition of four positions,
$\{\{1,3\},\{2,4\}\}$, has two directed spanning trees and
accounts for the excess $2^{-q}\rho$ in the fourth moment.

\fig{fig06_gram_trees}{Directed trees distinguish coherent Gram spectra
from the free limit. (a) Identifying equal-colored vertices of the directed
four-cycle for the crossing partition gives two directed spanning trees.
(b) At unit sampling density the fourth limiting moment is
$14+2^{-d}$; complex dimension two gives $57/4$. (c) Exact integer
cofactor enumeration gives the relative excess $M_p/C_p-1$ over Catalan moments at orders
four through seven. The excess decreases as complex dimension grows,
while higher moments retain more crossing contributions at fixed
dimension. The figure uses exact partition formulas, without sampled
eigenvalue data; $d$ in its axes is the complex dimension denoted $q$ in
the text.}{fig:gram}

\begin{lemma}[The local graph integral]
\label{lem:coherent-graph-integral}
Let $G$ be a fixed connected Eulerian directed multigraph on $r$
vertices, with multiplicities $n_{ab}$.  Then
\begin{equation}
 I_m(G):=\int\prod_{a,b}\langle z_a,z_b\rangle^{m n_{ab}}
                   \,d\sigma(z_1)\cdots d\sigma(z_r)
 =D_m^{-(r-1)}\tau(G)^{-q}(1+o(1)).
 \label{eq:coherent-graph-integral}
\end{equation}
Furthermore $|I_m(G)|\leq C_GD_m^{-(r-1)}$ for all sufficiently
large $m$.
\end{lemma}

\begin{proof}
The Eulerian condition makes the integrand independent of the lifts.
Unitary invariance permits fixing $z_1=[1:0:\cdots:0]$.  The
integrand has modulus at most one, with equality only when all rays
coincide.  Indeed equality forces the endpoints of every edge to
coincide projectively, and the graph is connected.  Away from a
fixed neighborhood of this point, the integral is exponentially
small in $m$.

For the remaining rays use the lifts
$(1,w_a)/(1+|w_a|^2)^{1/2}$, with $w_1=0$.  The integrand near
zero is $\exp(m\Psi(w))$, where
\[
 \Psi(w)=\sum_{a,b}n_{ab}\log(1+\overline w_a\cdot w_b)
       -\sum_a\left(\sum_b n_{ab}\right)\log(1+|w_a|^2).
\]
The Eulerian identity equates indegrees with outdegrees and accounts
for the second sum.  For each of the $q$ complex coordinate
directions the quadratic term is $-\overline w^{\,T}L^{(1)}w$.
The Hermitian part of $L^{(1)}$ is the positive definite pinned
Laplacian with undirected edge weights $(n_{ab}+n_{ba})/2$.
Consequently $\operatorname{Re}\Psi(w)\leq-c\sum_a|w_a|^2$
on a sufficiently small neighborhood.

The affine Fubini--Study density is
$(q!/\pi^q)(1+|w|^2)^{-(q+1)}d^{2q}w$.
After $w_a=m^{-1/2}\xi_a$, Gaussian domination gives
\[
 I_m(G)=(q!)^{r-1}m^{-q(r-1)}\det(L^{(1)})^{-q}(1+o(1)).
\]
Here we used
$\int_{\mathbb C^s}e^{-\overline\xi^{\,T}A\xi}d^{2s}\xi
=\pi^s/\det A$ when the Hermitian part of $A$ is positive
definite.  This identity follows by unitary Schur triangularization
and successive one-variable Gaussian integration.  Since
$D_m\sim m^q/q!$, the asserted asymptotic follows.  The same
Gaussian domination and the exponentially small exterior contribution
give the uniform bound.
\end{proof}

\begin{proof}[Proof of Theorem~\ref{thm:coherent-spectral-law}]
Expanding the trace according to equality patterns of particle
indices gives the exact identity
\begin{equation}
 \mathbb E\frac1n\operatorname{tr}\bigl[(B^{(m)})^p\bigr]
 =\sum_{\pi\in\mathfrak P(p)}
       \frac{(n)_{|\pi|}}n I_m(G_\pi).
 \label{eq:coherent-exact-trace-expansion}
\end{equation}
For a partition with $r$ blocks, Lemma~\ref{lem:coherent-graph-integral}
and $n/D_m\to\rho$ give its limit
$\rho^{r-1}\tau(G_\pi)^{-q}$.

For the variance, expand two trace words.  Disjoint label sets have
zero covariance.  When the words share at least one label, their
combined graph is connected and Eulerian.  If it has $r$ vertices,
its integral is $O(D_m^{-(r-1)})$ and there are $O(n^r)$ labelings.
The two trace normalizations leave $O(n^{-1})$.  Products of the
separate means have the same or smaller order: if the separate words
use $r_1,r_2$ vertices and share $s\geq1$, then
$r=r_1+r_2-s$, giving $O(n^{-s})$.  There are only finitely many
patterns for fixed $p$.

The spectral measures are tight since their first moments equal one.
The bounds on higher moments imply uniform integrability of every
lower moment.  Thus subsequential limits of the expected spectral
measures exist and have all the displayed moments.  Since
$\tau\geq1$,
\[
 M_{2p}(q,\rho)
 \leq \max(1,\rho^{2p-1})(2p)^{2p}.
\]
Carleman's criterion gives uniqueness of the moment problem.
The variance estimate then gives weak and moment convergence in
probability.  Chebyshev's inequality and Borel--Cantelli give the
almost-sure assertion whenever $\sum_m n_m^{-1}<\infty$; no
independence between different rows is required.

For the final assertion let $\Omega_m$ be the prescribed moment
event.  Boundedness and independence give
$\Pr(\Omega_m^c)=O((n_m\epsilon_m^2)^{-1})=o(1)$.
For $T_{m,p}=n_m^{-1}\operatorname{tr}\bigl[(B^{(m)})^p\bigr]$,
\[
 \mathbb E\bigl[|T_{m,p}-\mathbb ET_{m,p}|^2\mid\Omega_m\bigr]
 \leq\frac{\operatorname{Var}(T_{m,p})}{\Pr(\Omega_m)}
 =O_p(n_m^{-1}).
\]
This proves the same concentration, moment convergence, and
summability assertions for the conditioned rows.
\end{proof}

\subsection{Unbounded support and the emergence of noncrossing partitions}

The fixed-dimensional laws are not simply Marchenko--Pastur laws with
a shifted parameter.  Their support is unbounded, even though the
large-dimension limit has compact support.

\begin{lemma}[The noncrossing layer]
\label{lem:tree-one-noncrossing}
For a cycle quotient, $\tau(G_\pi)=1$ if and only if $\pi$ is
noncrossing in the cyclic order.
\end{lemma}

\begin{proof}
A noncrossing partition with more than one block has a proper interval
block.  Its quotient vertex has one nonloop incoming edge and one
nonloop outgoing edge.  Deleting this vertex and joining these edges
is a series reduction with Laplacian diagonal entry one; a Schur
complement shows that a cofactor rooted elsewhere is unchanged.
Induction ends with the one-vertex graph.

Conversely, a crossing supplies blocks $a,b$ occurring in order
$a,\ldots,b,\ldots,a,\ldots,b$.  At least two intervals between
consecutive occurrences of $a$ contain $b$.  Follow each interval
to its first $b$ and erase internal cycles.  This gives two simple
directed paths from $a$ to $b$ with distinct initial edges; parallel
edges are distinguished.  Each extends to an in-arborescence rooted
at $b$, by attaching remaining vertices along paths to the current
tree.  Such paths exist because a connected Eulerian graph is
strongly connected.  The two arborescences choose different outgoing
edges at $a$, so $\tau\geq2$.
\end{proof}

\begin{theorem}[Support, exponential moments, and the free limit]
\label{thm:coherent-free-limit}
For every finite positive integer $q$ and $\rho>0$, the support of
$\nu_{q,\rho}$ is unbounded.  Nevertheless its moment-generating
function is entire, and for $t\geq0$,
\begin{equation}
 \int e^{tx}\,d\nu_{q,\rho}(x)
 \leq1+\frac{\exp\{\rho(e^t-1)\}-1}{\rho}.
 \label{eq:coherent-entire-bound}
\end{equation}
In particular, for $R\geq\rho$,
\begin{equation}
 \nu_{q,\rho}([R,\infty))
 \leq\max(1,\rho^{-1})
       \exp\{-R\log(R/\rho)+R-\rho\}.
 \label{eq:coherent-poisson-tail}
\end{equation}
As $q\to\infty$ through integers, $\nu_{q,\rho}$ converges weakly
and in every fixed positive moment to the mean-one
Marchenko--Pastur law of aspect ratio $\rho$.
\end{theorem}

\begin{proof}
Consider all partitions of $2r$ positions into pairs.  Each quotient
vertex has outdegree at most two, after ignoring loops.  An
in-arborescence chooses one outgoing edge at every nonroot vertex,
so $\tau\leq2^{r-1}$.  There are $(2r-1)!!$ pair partitions.
Equation~\eqref{eq:coherent-tree-moments} therefore gives
\begin{equation}
 M_{2r}(q,\rho)
 \geq (2r-1)!!\left(\frac{\rho}{2^q}\right)^{r-1}.
 \label{eq:coherent-pairing-lower-bound}
\end{equation}
The $2r$-th root of the right side diverges with $r$.  A probability
measure with bounded support could not satisfy this inequality.

For the upper bound let
$T_p(\rho)=\sum_{\pi\in\mathfrak P(p)}\rho^{|\pi|}$
be the Touchard polynomial.  For $p\geq1$,
$M_p(q,\rho)\leq T_p(\rho)/\rho$.
The exponential generating function
$\sum_{p\geq0}T_p(\rho)t^p/p!=\exp\{\rho(e^t-1)\}$,
together with monotone convergence, proves
\eqref{eq:coherent-entire-bound}.  Finiteness for every $t\geq0$
also gives an entire complex moment-generating function.  The right
side of \eqref{eq:coherent-entire-bound} is at most
$\max(1,\rho^{-1})\exp\{\rho(e^t-1)\}$.
Markov's inequality with $t=\log(R/\rho)$ proves
\eqref{eq:coherent-poisson-tail}.

Finally, Lemma~\ref{lem:tree-one-noncrossing} gives
\[
 M_p(q,\rho)=
 \sum_{\pi\in NC(p)}\rho^{|\pi|-1}
 +\sum_{\pi\notin NC(p)}\rho^{|\pi|-1}\tau(G_\pi)^{-q}.
\]
The second sum is at most
$2^{-q}\sum_{\pi\notin NC(p)}\rho^{|\pi|-1}$.
The first is the $p$-th mean-one Marchenko--Pastur moment:
Kreweras complementation interchanges the exponents $|\pi|-1$
and $p-|\pi|$ in the usual noncrossing formula
\cite{prob_nica_speicher2006}.  The uniform Touchard bounds yield
tightness and uniform integrability of fixed moments.  Moment
determinacy of the compactly supported limiting law completes the
proof.
\end{proof}

\subsection{Finite-degree identities and the logarithmic scale}

The first moments admit simple exact checks before taking a limit.
With $D=D_m$ and $D'=D_{2m}$, the identities
$\int P_z^{\otimes m}d\sigma=I/D$ and
$\int|\langle z,w\rangle|^{4m}d\sigma(w)=1/D'$ imply
\begin{align}
 \mathbb E T_{m,2}&=1+\frac{n-1}{D},\notag\\
 \mathbb E T_{m,3}&=1+\frac{3(n-1)}D
                     +\frac{(n-1)(n-2)}{D^2},\notag\\
 \mathbb E T_{m,4}&=1+\frac{6(n-1)}D+\frac{n-1}{D'}
   +\frac{6(n-1)(n-2)}{D^2}
   +\frac{(n-1)(n-2)(n-3)}{D^3}.
 \label{eq:coherent-finite-moments}
\end{align}
The supplemental verifier enumerates cycle partitions, computes all
rooted cofactors by integer elimination, and checks the noncrossing
criterion and the rational moment identities.  These checks concern
finite combinatorics; numerical spectra provide illustrations of the
limiting law.

For anticanonical degree $m=(q+1)k$ and $n=D_m$, let $\Delta_k$
be the Slater determinant of a Fubini--Study orthonormal basis.
The evaluation kernel is $D_m\langle z,w\rangle^m$, so
\begin{equation}
 E_k:=-\frac1{kD_m}\log\|\Delta_k\|^2
 =-\frac1k\int\log\lambda\,d\nu_m(\lambda)
    -\frac{\log D_m}{k}.
 \label{eq:coherent-logarithmic-energy}
\end{equation}
The spectral law consequently gives a precise object for investigating
the collective determinant energy.  Its natural continuation is a
logarithmic-tail theorem on scales $\lambda\asymp e^{-kt}$,
and its analogue under attractive Gibbs reweighting.  Such a theorem
would connect coherent spectral statistics with the moment-reduced
large-deviation program for K\"ahler--Einstein measures
\cite{AndreassonBermanSvensson2026}.  Equation~\eqref{eq:coherent-logarithmic-energy}
also specifies the distinct scales involved: fixed positive eigenvalue
scales contribute only $o(1)$ to this normalized energy.
